%==============================================================================
\section{Data Structure --- Cấu trúc dữ liệu trong học máy}
%==============================================================================

\begin{dongco}[title={Vì sao cấu trúc dữ liệu quan trọng trong ML?}]
Cấu trúc dữ liệu đóng vai trò quan trọng trong học máy vì giúp tổ chức, thao tác và phân tích dữ liệu.
Dữ liệu là nền tảng của mô hình ML, và cấu trúc dữ liệu được dùng có thể ảnh hưởng đáng kể đến hiệu năng và độ chính xác.
\end{dongco}

\begin{ynghia}[title={target của bài}]
Cấu trúc dữ liệu giúp xây và hiểu các bài toán phức tạp trong ML.
Chọn cấu trúc dữ liệu cẩn thận có thể giúp tăng hiệu năng và tối ưu mô hình.
Bài này giới thiệu một số cấu trúc dữ liệu hay dùng và cách chúng được dùng trong ML.
\end{ynghia}

\subsection{Data Structure là gì?}

\begin{dinhnghia}[title={Định nghĩa}]
Cấu trúc dữ liệu là cách tổ chức và lưu trữ dữ liệu để sử dụng hiệu quả.
Chúng gồm các cấu trúc như array, linked list, stack, \ldots được thiết kế để hỗ trợ các thao tác cụ thể.
Trong ML, cấu trúc dữ liệu đặc biệt quan trọng ở tiền xử lý, cài đặt thuật toán và tối ưu hoá.
\end{dinhnghia}

\subsection{Các cấu trúc dữ liệu thường dùng trong ML}

\begin{dongco}[title={Chọn đúng cấu trúc giúp làm nhanh và tiết kiệm bộ nhớ}]
Cấu trúc dữ liệu phù hợp giúp xử lý nhanh hơn, truy cập dữ liệu dễ hơn và lưu trữ hiệu quả hơn.
\end{dongco}

\subsubsection{Arrays}
\begin{dinhnghia}[title={Arrays}]
Array là cấu trúc cơ bản để lưu và thao tác dữ liệu trong ML.
Phần tử truy cập qua chỉ số (index). Array cho truy xuất nhanh vì dữ liệu nằm liên tiếp trong bộ nhớ.
\end{dinhnghia}

\begin{ynghia}[title={Vì sao dùng array trong ML?}]
Vì có thể thực hiện phép toán vector hoá (vectorized operations) trên array, nên biểu diễn input data dạng array là lựa chọn tốt.
\end{ynghia}

\begin{phuongphap}[title={Một số tác vụ dùng array}]
\begin{itemize}
  \item Dữ liệu thô thường biểu diễn dưới dạng array.
  \item Chuyển pandas DataFrame sang list (lý do: pandas series yêu cầu cùng kiểu; Python list có thể chứa nhiều kiểu).
  \item Tiền xử lý: normalization, scaling, reshaping.
  \item Word embedding: tạo ma trận đa chiều.
\end{itemize}
\end{phuongphap}

\begin{luuy}[title={Hạn chế của array}]
Array dễ dùng và index nhanh, nhưng kích thước cố định, có thể là hạn chế khi làm với dataset lớn.
\end{luuy}

\subsubsection{Lists}
\begin{dinhnghia}[title={Lists}]
List là tập các phần tử có thể khác kiểu (heterogeneous), truy cập bằng iterator.
Trong ML, list thường dùng để lưu cấu trúc dữ liệu phức tạp như nested lists, dictionaries, tuples.
List linh hoạt và hỗ trợ kích thước thay đổi, nhưng thường chậm hơn array do cần iteration.
\end{dinhnghia}

\subsubsection{Dictionaries}
\begin{dinhnghia}[title={Dictionaries}]
Dictionary là tập các cặp key--value, truy cập theo key.
Thường dùng để lưu metadata hoặc label gắn với dữ liệu; hữu ích để tạo lookup table.
Dictionary truy cập nhanh nhưng có thể tốn bộ nhớ khi dataset lớn.
\end{dinhnghia}

\subsubsection{Linked Lists}
\begin{dinhnghia}[title={Linked lists}]
Linked list là tập các node, mỗi node chứa dữ liệu và tham chiếu tới node tiếp theo.
Trong ML, có thể dùng để lưu và thao tác dữ liệu tuần tự như time-series.
Linked list chèn/xoá hiệu quả nhưng truy cập dữ liệu thường chậm hơn array/list.
\end{dinhnghia}

\begin{luuy}[title={Nhận xét}]
Linked list hữu ích cho dữ liệu động (thêm/bớt nhiều), nhưng ít phổ biến hơn array vì array truy xuất dữ liệu hiệu quả hơn.
\end{luuy}

\subsubsection{Stack and Queue}

\begin{dinhnghia}[title={Stack (LIFO)}]
Stack theo LIFO (Last In First Out). Cách tiếp cận “stacking classifier” có thể dùng để giải multi-class bằng cách chia thành nhiều bài toán nhị phân, rồi “stack” các output để đưa vào meta-classifier.
\end{dinhnghia}

\begin{dinhnghia}[title={Queue (FIFO)}]
Queue theo FIFO (First In First Out), giống người xếp hàng.
Queue dùng trong multi-threading để tối ưu và điều phối luồng dữ liệu giữa các thread; thường dùng để xử lý dữ liệu lớn và feed batch cho training để train liên tục và hiệu quả.
\end{dinhnghia}

\subsubsection{Trees}
\begin{dinhnghia}[title={Trees}]
Tree là cấu trúc phân cấp, thường dùng trong thuật toán ra quyết định như decision trees và random forests.
Tree cho thuật toán tìm kiếm/sắp xếp hiệu quả nhưng có thể phức tạp khi cài đặt và có thể bị overfitting.
\end{dinhnghia}

\begin{luuy}[title={Binary trees}]
Binary trees như một dạng tree, cũng hay dùng trong decision trees và random forests, với các đặc tính tương tự.
\end{luuy}

\subsubsection{Graphs}
\begin{dinhnghia}[title={Graphs}]
Graph gồm nodes và edges, dùng để biểu diễn quan hệ phức tạp giữa các điểm dữ liệu.
Adjacency matrices và linked lists có thể dùng để tạo và thao tác graph.
Graph có thuật toán mạnh cho clustering, classification và prediction, nhưng có thể phức tạp và gặp vấn đề scalability.
\end{dinhnghia}

\begin{ynghia}[title={Ứng dụng graph}]
Graph dùng nhiều trong recommendation systems, link prediction và phân tích mạng xã hội.
\end{ynghia}

\subsubsection{Hash Maps}
\begin{dinhnghia}[title={Hash maps}]
Hash maps thường dùng trong ML nhờ khả năng lưu và truy xuất key--value.
Chúng thường dùng để lưu metadata/label; hữu ích để tạo lookup tables; nhưng có thể tốn bộ nhớ khi dataset lớn.
\end{dinhnghia}

\begin{luuy}[title={Cấu trúc chuyên biệt}]
Ngoài các cấu trúc trên, nhiều thư viện/framework ML còn cung cấp cấu trúc chuyên biệt cho use case cụ thể, như matrices và tensors cho deep learning.
Khi chọn cấu trúc dữ liệu, nên cân nhắc: kích thước dữ liệu, tốc độ xử lý và bộ nhớ.
\end{luuy}

\subsection{Cấu trúc dữ liệu được dùng trong ML như thế nào?}

\subsubsection{Storing and Accessing Data}
\begin{phuongphap}[title={Lưu và truy cập dữ liệu}]
Thuật toán ML cần nhiều dữ liệu cho train/test.
Arrays, lists, dictionaries giúp lưu và truy cập hiệu quả.
Ví dụ: array lưu tập số; dictionary lưu metadata/label gắn với dữ liệu.
\end{phuongphap}

\subsubsection{Pre-processing Data}
\begin{phuongphap}[title={Tiền xử lý dữ liệu}]
Trước khi train mô hình, cần làm sạch, biến đổi và chuẩn hoá.
Lists và arrays có thể dùng để lưu và thao tác dữ liệu trong tiền xử lý.
Ví dụ: list lọc missing values; array chuẩn hoá dữ liệu.
\end{phuongphap}

\subsubsection{Creating Feature Vectors}
\begin{phuongphap}[title={Tạo vector đặc trưng}]
Vector đặc trưng là thành phần quan trọng vì biểu diễn feature dùng để dự đoán.
Arrays và matrices thường dùng để tạo feature vectors.
Ví dụ: array lưu pixel ảnh; matrix lưu phân phối tần suất từ trong tài liệu văn bản.
\end{phuongphap}

\subsubsection{Building Decision Trees}
\begin{phuongphap}[title={Xây cây quyết định}]
Cây quyết định dùng cấu trúc tree để ra quyết định dựa trên tập feature.
Cây quyết định hữu ích cho classification và regression.
Cây được tạo bằng cách chia dữ liệu đệ quy theo feature “thông tin nhất”; cấu trúc tree giúp dễ duyệt quá trình ra quyết định và dự đoán.
\end{phuongphap}

\subsubsection{Building Graphs}
\begin{phuongphap}[title={Xây dựng graph}]
Graph dùng để biểu diễn quan hệ phức tạp giữa các điểm dữ liệu.
Adjacency matrices và linked lists dùng để tạo/thao tác graph.
Graph dùng cho clustering, classification và prediction.
\end{phuongphap}